\documentclass[11pt]{article}

\usepackage[utf8]{inputenc}
\usepackage[T1]{fontenc}
\usepackage{geometry}
\usepackage{amsmath}
\usepackage{amssymb}
\usepackage{booktabs}
\usepackage{hyperref}
\usepackage{xurl}
\usepackage{xcolor}
\usepackage{listings}
\usepackage{enumitem}
\usepackage{parskip}
\usepackage{graphicx}
\graphicspath{{../images/}}

\geometry{margin=1in}

\lstset{
  language=Java,
  basicstyle=\ttfamily\small,
  keywordstyle=\color{blue},
  commentstyle=\color{gray},
  stringstyle=\color{teal},
  showstringspaces=false,
  columns=fullflexible,
  frame=single,
  framerule=0.2pt,
  breaklines=true
}

\setlist[itemize]{topsep=0.3em,itemsep=0.2em}
\setlist[enumerate]{topsep=0.3em,itemsep=0.2em}

% Simple per-section source line.
\newcommand{\notesources}[1]{\par\noindent{\small\color{gray}\textit{Sources:} \sloppy #1\par}}

% Unit-wise source strings for this subject (used by slides/notes).
% Path is relative to the lecture `latex/` folder (the compilation CWD).
% OOPS with Java (CSEG1044) — unit-wise sources
%
% These are short, student-facing reference strings intended to be used with:
%   - \slidesources{...}  (in beamer slides)
%   - \notesources{...}   (in lecture notes)
%
% Style inspiration: other subjects in My-Subjects/ use semicolon-separated sources
% and include an "accessed" date for web documentation.

\newcommand{\OOPSUnitISources}{Schildt, \textit{Java: A Beginner's Guide} (10e), McGraw-Hill (Java fundamentals + OOP basics).; Oracle Java Tutorials: Getting Started + Language Basics + Classes and Objects (accessed \today).; Java SE API docs (e.g., \texttt{String}, \texttt{Scanner}) (accessed \today).}

\newcommand{\OOPSUnitIISources}{Schildt, \textit{Java: The Complete Reference} (12e), McGraw-Hill (classes, inheritance, packages, interfaces).; Bloch, \textit{Effective Java} (3e), Addison-Wesley, 2018 (design choices; interfaces vs abstract classes).; Oracle Java Tutorials: Classes and Objects + Interfaces + Packages (accessed \today).; Java SE API docs (e.g., \texttt{Comparable}, \texttt{Runnable}) (accessed \today).}

\newcommand{\OOPSUnitIIISources}{Schildt, \textit{Java: The Complete Reference} (12e) (nested/inner classes, exceptions, threads, I/O).; Oracle Java Tutorials: Nested Classes + Exceptions + Concurrency + Basic I/O (accessed \today).; Java SE API docs (e.g., \texttt{Thread}, \texttt{AutoCloseable}) (accessed \today).}

\newcommand{\OOPSUnitIVSources}{Schildt, \textit{Java: The Complete Reference} (12e) (generics, lambdas, Swing, JDBC).; Bloch, \textit{Effective Java} (3e) (generics + lambdas).; Oracle Java Tutorials: Generics + Lambda Expressions + Swing + JDBC Basics (accessed \today).; Java SE API docs (e.g., \texttt{java.util.function}, \texttt{javax.swing}, \texttt{java.sql}) (accessed \today).}

\newcommand{\OOPSUnitVSources}{Schildt, \textit{Java: The Complete Reference} (12e) (Collections Framework + wrapper classes).; Oracle Java docs: Collections Framework overview (accessed \today).; Java SE API docs (\texttt{java.util} collections; wrapper classes in \texttt{java.lang}) (accessed \today).}

\newcommand{\OOPSUnitVISources}{Git SCM: \textit{Pro Git} (2e) (version control workflow), accessed \today.; GitHub Docs (repositories, issues, pull requests), accessed \today.; Oracle Java Tutorials: Swing + JDBC (accessed \today).; JUnit 5 User Guide (accessed \today).; Apache Maven Project: Getting Started (accessed \today).; Gradle User Manual: Getting Started (accessed \today).; UML class diagrams: UML 2.x overview/spec (accessed \today).}

\newcommand{\OOPSMiscWhyInterfacesSources}{Schildt, \textit{Java: The Complete Reference} (12e) (interfaces).; Bloch, \textit{Effective Java} (3e) (prefer interfaces where appropriate).; Oracle Java Tutorials: Interfaces (accessed \today).; Java SE API docs (e.g., \texttt{Comparable}, \texttt{Runnable}) (accessed \today).}



\title{Object Oriented Programming (CSEG1044)\\Unit 04 -- Lecture 04 Notes\\Swing Basics (Components, Containers, Layouts)}
\begin{document}
\maketitle
\tableofcontents

\section{Lecture Snapshot}
\begin{itemize}
  \item \textbf{Unit focus:} Generics, Lambdas, Swing, and JDBC
  \item \textbf{Main new ideas:} Swing overview: \texttt{JFrame}, \texttt{JPanel}, components vs containers, Layout managers (Flow/Border/Grid), UI composition: multi-panel skeleton
  \item \textbf{Course thread:} Use Swing Basics (Components, Containers, Layouts) to move Campus Resource Manager toward a real desktop application with UI, callbacks, and persistence.
\end{itemize}
\notesources{\OOPSUnitIVSources}

\section{Lecture Overview}
Swing is Java's GUI toolkit for building desktop applications.
In this lecture we learn the minimum you need to build a clean UI skeleton:
\begin{itemize}
  \item top-level window (\texttt{JFrame}),
  \item containers (\texttt{JPanel}) and components (\texttt{JLabel}, \texttt{JButton}, \texttt{JTextField}),
  \item layout managers (Flow/Border/Grid),
  \item composing multiple panels to create a complete screen.
\end{itemize}
\notesources{\OOPSUnitIVSources}

\section{Prerequisite Bridge}
Before reading the new material in depth, do a fast recall of the older ideas listed below.
\begin{itemize}
  \item \textbf{Classes and packages}: Swing applications still depend on clear classes and package boundaries. Main reference: \texttt{Unit 02 / Lecture 03 slides/notes}.
  \item \textbf{Nested helpers}: GUI code often uses callback and helper patterns introduced earlier. Main reference: \texttt{Unit 03 / Lecture 01 slides/notes}.
\end{itemize}
This lecture only revisits those ideas briefly so that the main explanation can stay focused on the new concept sequence.
\notesources{\OOPSUnitIVSources}

\section{Learning Outcomes}
After this lecture, you should be able to:
\begin{itemize}
  \item Explain Swing component model: \texttt{JFrame}, \texttt{JPanel}, components vs containers.
  \item Use layout managers (Flow/Border/Grid) to structure a UI.
  \item Build and run a small multi-panel UI skeleton.
\end{itemize}

\section{Suggested Reading Order}
\begin{enumerate}
  \item Read the \textbf{Prerequisite Bridge} first and confirm each recalled idea.
  \item Study the central explanation in this order: \textit{Swing overview: \texttt{JFrame}, \texttt{JPanel}, components vs containers} $\rightarrow$ \textit{Layout managers (Flow/Border/Grid)} $\rightarrow$ \textit{UI composition: multi-panel skeleton}.
  \item Trace the worked example line by line before checking short answers.
  \item Attempt the practice section without looking at the solution immediately.
  \item Finish with the exit questions and further-study suggestions to confirm retention.
\end{enumerate}
\notesources{\OOPSUnitIVSources}

\section{Key Terms (Quick)}
\begin{itemize}
  \item \textbf{Component}: UI element like button, label, text field.
  \item \textbf{Container}: holds components (e.g., \texttt{JPanel}, \texttt{JFrame} content pane).
  \item \textbf{Layout manager}: decides how components are arranged (sizes/positions).
  \item \textbf{Event Dispatch Thread (EDT)}: Swing UI thread (advanced but important for correctness).
  \item \textbf{MVC idea}: separate UI from business logic (we will use this later with JDBC).
\end{itemize}

\section{Detailed Notes}
\subsection{Swing overview: \texttt{JFrame}, \texttt{JPanel}, components vs containers}
Swing applications usually have:
\begin{itemize}
  \item a top-level window: \texttt{JFrame},
  \item inside it, one or more \texttt{JPanel} containers,
  \item inside panels, UI components: \texttt{JLabel}, \texttt{JTextField}, \texttt{JButton}, etc.
\end{itemize}

\textbf{Best practice:}
\begin{itemize}
  \item Build UI first (layout + widgets), then add event handling (next lecture).
  \item Keep UI code separate from core logic (helps testing and maintenance).
  \item Call \texttt{setVisible(true)} at the end, after adding components.
\end{itemize}

\subsection{Layout managers (Flow/Border/Grid)}
Swing layouts prevent hard-coded pixel positions and make UI more robust.

\textbf{FlowLayout:}
\begin{itemize}
  \item Places components left-to-right, wrapping to next line.
  \item Good for small button rows.
\end{itemize}

\textbf{BorderLayout:}
\begin{itemize}
  \item 5 regions: \texttt{NORTH}, \texttt{SOUTH}, \texttt{EAST}, \texttt{WEST}, \texttt{CENTER}.
  \item Great for ``header + main content + footer'' screens.
\end{itemize}

\textbf{GridLayout:}
\begin{itemize}
  \item Uniform rows/columns.
  \item Good for forms (label + field pairs) and simple dashboards.
\end{itemize}

\subsection{UI composition: multi-panel skeleton}
A very common pattern for beginners:
\begin{itemize}
  \item \textbf{Root panel}: \texttt{BorderLayout}
  \item \textbf{Top panel (NORTH)}: title label
  \item \textbf{Center panel}: form or table
  \item \textbf{Bottom panel (SOUTH)}: action buttons (save, clear, exit)
\end{itemize}
\notesources{\OOPSUnitIVSources}

\section{Running Project Connection}
Use Swing Basics (Components, Containers, Layouts) to move Campus Resource Manager toward a real desktop application with UI, callbacks, and persistence.
\begin{itemize}
  \item Use Campus Resource Manager as the running case study, or map the same idea to your own project domain if you are using another example.
  \item Ask where today's idea should live: model, service, DAO, UI, or team workflow.
  \item Use the lecture example as a smaller version of the same design choice you will make in the capstone.
\end{itemize}
\notesources{\OOPSUnitIVSources}

\section{Common Mistakes and Confusions}
\begin{itemize}
  \item Putting every control directly on one JFrame without structure.
  \item Hardcoding coordinates instead of using layout managers intentionally.
  \item Mixing UI code with business logic too early.
\end{itemize}
\notesources{\OOPSUnitIVSources}

\section{Worked Example: a multi-panel form skeleton}
\begin{lstlisting}
import javax.swing.*;
import java.awt.*;

public class SimpleFormUI {
    public static void main(String[] args) {
        SwingUtilities.invokeLater(() -> {
            JFrame f = new JFrame("Campus Resource Manager (Skeleton)");
            f.setDefaultCloseOperation(JFrame.EXIT_ON_CLOSE);

            JPanel root = new JPanel(new BorderLayout(10, 10));
            root.setBorder(BorderFactory.createEmptyBorder(10, 10, 10, 10));

            JLabel title = new JLabel("Add Resource", SwingConstants.CENTER);
            title.setFont(title.getFont().deriveFont(Font.BOLD, 18f));
            root.add(title, BorderLayout.NORTH);

            JPanel form = new JPanel(new GridLayout(3, 2, 8, 8));
            JTextField nameField = new JTextField();
            JTextField typeField = new JTextField();
            JTextField capacityField = new JTextField();
            form.add(new JLabel("Name:"));     form.add(nameField);
            form.add(new JLabel("Type:"));     form.add(typeField);
            form.add(new JLabel("Capacity:")); form.add(capacityField);
            root.add(form, BorderLayout.CENTER);

            JPanel actions = new JPanel(new FlowLayout(FlowLayout.RIGHT));
            JButton saveBtn = new JButton("Save");
            JButton clearBtn = new JButton("Clear");
            JButton exitBtn = new JButton("Exit");
            actions.add(saveBtn);
            actions.add(clearBtn);
            actions.add(exitBtn);
            root.add(actions, BorderLayout.SOUTH);

            f.setContentPane(root);
            f.pack();
            f.setLocationRelativeTo(null); // center on screen
            f.setVisible(true);
        });
    }
}
\end{lstlisting}

\section{Demo / Observation Task}
Use this block as a short reproducible demo or observation task.
\begin{itemize}
  \item Reproduce the lecture's main example around \textit{Swing overview: \texttt{JFrame}, \texttt{JPanel}, components vs containers} once without looking at the solution first.
  \item Change one input, rule, or design choice in Campus Resource Manager and predict the result before checking it.
  \item Record one success case, one edge case, and one failure case so the idea becomes observable instead of purely theoretical.
\end{itemize}
\notesources{\OOPSUnitIVSources}

\section{Practice Check (with brief solutions)}
\begin{enumerate}
  \item What is the difference between a component and a container?\par
  \textbf{Answer:} a component is a UI element; a container holds components.

  \item Which layout is good for header + main content + footer screens?\par
  \textbf{Answer:} \texttt{BorderLayout}.

  \item Why do we use layout managers instead of fixed pixel positions?\par
  \textbf{Answer:} layouts adapt better to different screen sizes/fonts and reduce UI bugs.
\end{enumerate}

\section{Self-Study Practice (no solutions)}
\begin{enumerate}
  \item Add one more field to the form: \texttt{Location}.
  \item Replace \texttt{GridLayout} with \texttt{GridBagLayout} (optional advanced) and observe flexibility.
  \item Create a second screen panel (e.g., \texttt{List Resources}) and switch panels (hint: \texttt{CardLayout}).
\end{enumerate}

\section{Further Study}
\begin{itemize}
  \item Revisit the prerequisite references if any recall bullet still feels weak.
  \item Rework the lecture example with one changed assumption, input, or design choice.
  \item Read the textbook and documentation references listed in the source line for a second explanation of the same topic.
  \item Apply the same idea once inside Campus Resource Manager so that the concept moves from theory to design practice.
\end{itemize}
\notesources{\OOPSUnitIVSources}

\section{Exit Questions (with sample answers)}
\textbf{Q1.} Why is \texttt{pack()} often preferred over \texttt{setSize()}?\par
\textbf{Sample answer:} \texttt{pack()} sizes the window based on preferred sizes of components and layout, producing a more consistent UI across systems.

\vspace{0.6em}
\textbf{Q2.} What is the role of \texttt{JPanel}?\par
\textbf{Sample answer:} \texttt{JPanel} is a container used to group and layout components. It helps structure the UI into manageable parts.

\vspace{0.6em}
\textbf{Q3.} Why do layout managers matter more than pixel coordinates in most Swing code?\par
\textbf{Sample answer:} Layouts adapt better to resizing, platform differences, and maintenance.

\end{document}
